\documentclass[11pt]{article}
\usepackage[utf8]{inputenc}
\usepackage{amsmath}
\usepackage{amssymb}
\usepackage{enumitem}
\usepackage{hyperref}
\usepackage{geometry}
\geometry{margin=1in}

\title{Vidhoor Legal Copilot - Backend Technology FAQ}
\author{}
\date{}

\begin{document}

\maketitle

\begin{enumerate}
\item \textbf{Question:} What is the primary web framework used in the Vidhoor Legal Copilot backend, and what version is specified in requirements.txt?
\\ \textbf{Answer:} The primary web framework is FastAPI, version 0.109.2 as specified in \texttt{backend/requirements.txt}.

\item \textbf{Question:} How does the application handle Cross-Origin Resource Sharing (CORS) to enable communication with the frontend?
\\ \textbf{Answer:} The application configures CORS middleware in \texttt{backend/main.py} (lines 33-46) to allow specific origins including localhost ports 3000, 5173, and 8080, with credentials enabled and all methods/headers allowed.

\item \textbf{Question:} Which LLM provider and model is primarily used for generating legal responses, and how is fallback implemented?
\\ \textbf{Answer:} The system uses Cerebras LLM with the primary model "llama3.1-8b". Fallback implementation is in \texttt{backend/llm\_engine.py} where \texttt{\_build\_model\_candidates} creates a list including "llama3.1-8b" and "gpt-oss-120b", and the \texttt{\_switch\_model} method tries alternatives if the primary fails.

\item \textbf{Question:} What vector database is used for storing and retrieving legal embeddings, and how is it accessed?
\\ \textbf{Answer:} ChromaDB is used as the vector database. It's accessed via HTTP client in \texttt{backend/chroma\_manager.py} (lines 286-292) connecting to a Chroma server at configurable host/port (default localhost:8000).

\item \textbf{Question:} How does the system implement hybrid retrieval combining vector search with BM25 ranking?
\\ \textbf{Answer:} Hybrid retrieval combines vector search (ChromaDB) with BM25 ranking from Oracle persistence. Weights are defined as HYBRID\_VECTOR\_WEIGHT = 0.5 and HYBRID\_BM25\_WEIGHT = 0.5 in \texttt{backend/chroma\_manager.py} (lines 38-39). The BM25 index is warmed from Oracle chunks and used alongside vector search.

\item \textbf{Question:} What Oracle database features are utilized for persisting chat history and evidence?
\\ \textbf{Answer:} The system uses Oracle Autonomous Database with:
\begin{itemize}
\item Tables for chat sessions (\texttt{vidhoor\_chat\_sessions}) and messages (\texttt{vidhoor\_chat\_messages})
\item Tables for evidence storage (\texttt{vidhoor\_user\_evidence})
\item CLOB fields for storing large text content
\item MERGE statements for upsert operations
\item Connection pooling with Oracle Wallet support
\item Schema initialization with backward compatibility
\end{itemize}

\item \textbf{Question:} How are PDF documents processed for text extraction in the OCR pipeline?
\\ \textbf{Answer:} PDF processing uses PyPDF (specified in requirements.txt) and is handled by \texttt{VisionOCRService} in \texttt{backend/services/ocr\_vision.py}. The service extracts text from PDF pages and processes them through OCR and translation pipelines.

\item \textbf{Question:} What PII (Personally Identifiable Information) protection mechanisms are implemented in the system?
\\ \textbf{Answer:} PII protection is implemented via:
\begin{itemize}
\item \texttt{PIIVault} class in \texttt{backend/pii\_vault.py} using Microsoft Presidio
\item Entity detection and anonymization before LLM processing
\item Masking of sensitive information in chat messages and evidence
\item Separate storage of PII mappings for potential unmasking by frontend
\end{itemize}

\item \textbf{Question:} How does the application handle multilingual support, particularly for translating non-English legal documents?
\\ \textbf{Answer:} Multilingual support uses:
\begin{itemize}
\item \texttt{langdetect} for language detection
\item Helsinki NLP translation model via \texttt{services/translate\_helsinki.py}
\item Translation of OCR-extracted text to English for processing
\item Language detection in OCR results (\texttt{OCRPageResult.detected\_language})
\end{itemize}

\item \textbf{Question:} What embedding models are used for converting legal text to vectors, and what fallback strategy exists?
\\ \textbf{Answer:} Primary embedding model is BAAI/bge-m3 with fallback to all-MiniLM-L6-v2. The fallback strategy is implemented in \texttt{backend/chroma\_manager.py} \texttt{\_build\_embedding\_function} method (lines 312-338) which tries the primary model first, then falls back on failure.

\item \textbf{Question:} How are legal citations structured and validated in the system?
\\ \textbf{Answer:} Legal citations are structured using the \texttt{Citation} Pydantic model (lines 60-69 in \texttt{backend/main.py}) with fields for doc\_id, title, source, source\_url, section, page, snippet, confidence, and last\_updated. Validation occurs through reference matching functions like \texttt{\_references\_match} and \texttt{\_reference\_match\_for\_fir}.

\item \textbf{Question:} What security measures are in place for handling encrypted evidence files?
\\ \textbf{Answer:} Evidence security includes:
\begin{itemize}
\item Encryption using AES-GCM (via \texttt{evidenceCrypto.ts} frontend, but backend handles storage)
\item Storage of encrypted payloads, IVs, and metadata in Oracle DB
\item Key management through key\_id references
\item Separation of encrypted content from masked summaries/analyses
\item Access control via user\_id and session\_id foreign keys
\end{itemize}

\item \textbf{Question:} How does the system determine if a user query is related to legal matters?
\\ \textbf{Answer:} Legal query detection uses:
\begin{itemize}
\item Keyword matching against \texttt{LEGAL\_QUERY\_KEYWORDS} set (lines 149-177 in \texttt{backend/main.py})
\item Regex pattern matching for article/section references (line 365-366)
\item Implemented in \texttt{is\_legal\_query()} function (lines 355-368)
\end{itemize}

\item \textbf{Question:} What is the purpose of the BM25Okapi implementation and how is it integrated with ChromaDB?
\\ \textbf{Answer:} BM25Okapi provides lexical ranking to complement semantic vector search. Integration occurs through:
\begin{itemize}
\item Persistence of chunks to Oracle via \texttt{OracleChunkRepository}
\item In-memory BM25 index built from Oracle chunks (\texttt{refresh\_bm25\_from\_oracle})
\item Hybrid scoring combining vector distance and BM25 scores in \texttt{retrieve\_context}
\end{itemize}

\item \textbf{Question:} How are legal acts (like BNS, BNSS, BSA) identified and filtered during retrieval?
\\ \textbf{Answer:} Legal act identification uses:
\begin{itemize}
\item \texttt{infer\_act\_filter()} and \texttt{infer\_act\_filters()} functions (lines 297-352)
\item Pattern matching for act names and abbreviations (BNS, BNSS, BSA)
\item Special handling for bail-related queries to default to BNSS
\item Filtering in \texttt{\_source\_matches\_act\_filter()} to prevent cross-contamination
\end{itemize}

\item \textbf{Question:} What role does the PII Vault play in the evidence handling workflow?
\\ \textbf{Answer:} The PII Vault:
\begin{itemize}
\item Detects and masks PII in evidence before storage
\item Maintains mapping between original and masked values
\item Provides unmasking capabilities for authorized frontend display
\item Is invoked during OCR processing and evidence saving workflows
\end{itemize}

\item \textbf{Question:} How does the system manage session persistence and chat history in Oracle Database?
\\ \textbf{Answer:} Session management uses:
\begin{itemize}
\item \texttt{vidhoor\_chat\_sessions} table for session metadata
\item \texttt{vidhoor\_chat\_messages} table for message history with CLOB content
\item Oracle MERGE operations for session creation/update
\item Timestamps for tracking creation and updates
\item Foreign key constraints maintaining session-message relationships
\item Pinned session feature for user preferences
\end{itemize}

\item \textbf{Question:} What is the function of the \texttt{infer\_act\_filter} and \texttt{infer\_act\_filters} functions?
\\ \textbf{Answer:} These functions analyze user queries to determine relevant legal acts for retrieval filtering:
\begin{itemize}
\item \texttt{infer\_act\_filter()} returns single most likely act
\item \texttt{infer\_act\_filters()} returns list of possible acts
\item Both use keyword and pattern matching to identify BNS, BNSS, BSA, Constitution
\item Special handling for bail queries to prioritize BNSS
\end{itemize}

\item \textbf{Question:} How does the system handle OCR processing for legal documents, including language detection?
\\ \textbf{Answer:} OCR processing flow:
\begin{enumerate}
\item PDF/image input processed by \texttt{VisionOCRService}
\item Text extraction per page with language detection (\texttt{langdetect})
\item Non-English text translated to English via Helsinki model
\item Masking of PII in extracted text
\item Storage of both original and translated/masked results
\item Language tracking in \texttt{OCRPageResult.detected\_language}
\end{enumerate}

\item \textbf{Question:} What is the purpose of the \texttt{offence\_guidance\_rules} and how are they applied?
\\ \textbf{Answer:} Offence guidance rules provide practical legal guidance when strict citation grounding isn't possible:
\begin{itemize}
\item Defined in \texttt{OFFENCE\_GUIDANCE\_RULES} (lines 189-210)
\item Contain patterns, applicable acts, section hints, and guidance bullet points
\item Applied via \texttt{\_detect\_offence\_rule()} and \texttt{\_build\_offence\_guidance\_markdown()}
\item Used in OCR analysis when citation confidence is low
\end{itemize}

\item \textbf{Question:} How are legal chunks ingested into the system, and what metadata is associated with them?
\\ \textbf{Answer:} Ingestion process:
\begin{enumerate}
\item Legal texts chunked via \texttt{ingest\_legal\_resources.py} or \texttt{ingest\_constitution.py}
\item Metadata includes: source, act, section/article, status
\item Stored in ChromaDB with vector embeddings
\item Persisted to Oracle for BM25 recovery via \texttt{\_persist\_chunks\_for\_bm25}
\item Chunk IDs generated deterministically from content hash
\end{enumerate}

\item \textbf{Question:} What is the role of the \texttt{OracleChunkRepository} in the Chroma manager?
\\ \textbf{Answer:} The \texttt{OracleChunkRepository}:
\begin{itemize}
\item Persists legal chunks to Oracle Database for BM25 recovery
\item Provides \texttt{load\_chunks()} and \texttt{upsert\_chunks()} methods
\item Initializes schema for chunk storage table
\item Enables rebuilding BM25 index from persistent storage
\item Acts as backup when ChromaDB needs reconstruction
\end{itemize}

\item \textbf{Question:} How does the system ensure that LLM responses are strictly grounded in retrieved legal context?
\\ \textbf{Answer:} Response grounding is ensured by:
\begin{itemize}
\item Strict system prompt in \texttt{LLMEngine.prompt} (lines 78-116)
\item Instruction to ONLY answer using provided legal context
\item Requirement to say context is insufficient if answer not present
\item Prohibition against hallucination or inferring beyond exact matches
\item Structured response format with specific sections and bullet points
\item Post-processing to enforce bullet points under subheadings
\end{itemize}

\item \textbf{Question:} What environment variables are required for configuring the Oracle Autonomous Database connection?
\\ \textbf{Answer:} Required Oracle environment variables:
\begin{itemize}
\item \texttt{ORACLE\_USER} - database username
\item \texttt{ORACLE\_PASSWORD} - database password
\item \texttt{ORACLE\_DSN} - database connection string
\item Optional: \texttt{ORACLE\_CONFIG\_DIR}, \texttt{ORACLE\_WALLET\_LOCATION}, \texttt{ORACLE\_WALLET\_PASSWORD}
\end{itemize}
Used in \texttt{OracleChatHistoryRepository.\_\_init\_\_()} (lines 22-36).

\item \textbf{Question:} How does the system handle model switching and fallback when the primary Cerebras model is unavailable?
\\ \textbf{Answer:} Model fallback mechanism:
\begin{itemize}
\item \texttt{\_build\_model\_candidates()} creates ordered list: [primary\_model, "llama3.1-8b", "gpt-oss-120b"]
\item \texttt{\_switch\_model()} attempts to initialize alternative models
\item Generation methods (\texttt{generate\_legal\_response}, etc.) iterate through candidates
\item On model\_not\_found errors, logs warning and tries next candidate
\item Falls back to seed title generation if all models fail
\end{itemize}
\end{enumerate}

\end{document}